% ============================================================================
% 00_abstract.tex — auction-v3 (LLM-based mechanism-design screening reframe)
% \begin{abstract}...\end{abstract} only; no citations (CONVENTIONS §1).
% Per revision directive: non-technical opening for a general management
% audience; ordinal validation named as the main result with the evidence
% spelled out; ordinal-vs-cardinal folded into one takeaway; "stated
% reasoning" defined with an example; NO eBay, NO DA, NO third price,
% NO frontier material in the abstract.
% ============================================================================
\begin{abstract}
Well-designed market rules often fail in practice for a simple reason: participants do not play the strategy the rules were designed around. In a second-price auction the winner pays the runner-up's bid, so bidding exactly what the item is worth to you is optimal no matter what anyone else does---yet decades of laboratory experiments find that people persistently bid above their value, while the same rule played as an ascending price clock is followed almost correctly. Which implementations, explanations, and decision aids induce correct play is therefore an empirical design question, and the standard instrument for answering it---incentivized human experiments---is too slow and expensive to search a large design space. This paper evaluates a cheaper instrument: a panel of large language model (LLM) bidding agents, used to \emph{screen} candidate designs for which ones preserve correct play, not to simulate human bids. We report two results. First, an ordinal validation. On every auction comparison for which published human evidence exists---first-price versus second-price sealed bidding (harder versus easier), sealed bids versus ascending clocks (the clock helps), and menu-style restatements of the rules (no reliable effect)---a panel of four LLM families recovers the direction of the human finding, family by family, even though the two populations err in opposite directions (humans overbid; the models underbid) and no tuning of the agents reproduces human bidding in levels. The agreement is ordinal, not cardinal---and ordinal is exactly the signal design screening needs. Second, theory-guided discovery. We organize the design space along three dimensions from the theory of simple mechanisms---how a mechanism is implemented, how it is described, and what reasoning support participants receive---and use the panel to populate the resulting design map. Cells with human anchors reproduce the human record; cells without become concrete, falsifiable predictions for human experiments. The sharpest: one honest sentence stating \emph{why} truthful bidding is safe (``your bid determines whether you win, not what you pay'') recovers much of the clock's benefit in every model family, while restating the rules without that content does nothing on average, and prompting agents to reason about opponents backfires. Because each simulated bid is accompanied by a stated plan---a short written explanation of the intended strategy, such as ``I will bid slightly below my value to keep a profit margin''---we can also show that the effective designs change bids without changing the plans agents articulate: a screening instrument should audit behavior, not stated understanding.
\end{abstract}
